% ============================================================
% Introduction
% ============================================================
\section{Introduction}
\label{sec:intro}

As machine learning (ML) has advanced, identifying patterns in data has
increasingly shifted from a manual task to one delegated to learned
models. Both the volume of data and the range of data types have grown.
Labeling such data, however, is costly, so most of it remains unlabeled.
Demand has therefore risen for unsupervised learning, which extracts
meaningful structure without explicit labels.

Among such methods, clustering is used as an analysis tool in its own
right, for example in geographic information
systems~\cite{tu2022dbscan}, in outlier detection over
financial transactions~\cite{yang2014dbscan,ahmed2016survey}, and in the
analysis of healthcare data~\cite{khanizadeh2025smart,ahlqvist2018novel}.
It also serves as a building block for downstream pipelines such as
retrieval-augmented generation~\cite{sarthi2024raptor,edge2024local}.
In practice, clustering frequently operates on data aggregated from
multiple data owners or on directly sensitive records. Data privacy is
therefore a primary concern for clustering.

K-means is the most commonly studied algorithm in the privacy-preserving
clustering literature. Its update rule of additions and averages is
inherently friendly to homomorphic encryption (HE)~\cite{gentry2009fully},
and comparison is needed only to assign each point to its nearest
centroid. K-means, however, detects only convex clusters, cannot identify
noise, and requires the number of clusters to be fixed in advance---a
value that is rarely known in practice. This paper therefore targets Density-Based Spatial Clustering
of Applications with Noise (DBSCAN)~\cite{ester1996density}, which
discovers clusters of arbitrary shape, labels noise explicitly, and
requires no prior cluster count.

Executing DBSCAN over encrypted data, however, faces a single fundamental
barrier. DBSCAN is built on comparisons---whether a distance lies within
$\varepsilon$, whether a neighbor count is at least $\mathit{MinPts}$---and
homomorphic encryption provides no comparison as a native operation.
Existing schemes have circumvented this limitation in different ways.

The first approach returns the comparison, or other steps of the
computation, to the client~\cite{wang2022homomorphic,rahman2017towards}.
In the scheme of Wang et al.~\cite{wang2022homomorphic}, the server sends
the encrypted, shifted difference of the two compared values to the
client, which decrypts it and replies with the outcome of the comparison
as a single bit. This
recovers the missing primitive but introduces three problems. First, one
round trip is incurred per comparison, so the clustering is serialized on
client responses and its latency is bound to how often the client
replies. Second, each reply is a plaintext
bit, and the server knows which pair it queried; the accumulated replies
therefore yield the $\varepsilon$-neighborhood relation, the core, border,
and noise labels derived from it, and ultimately the cluster membership
itself, since the server executes the clustering control flow directly on
these plaintext outcomes. The coordinates are protected, but the structure
is not. Third, the security analysis of~\cite{wang2022homomorphic} argues
that the server cannot recover the data without breaking the encryption,
and considers the returned bits only as a possible aid to recovering the
key. It does not consider what these bits reveal about the structure of
the data.

The second approach adopts secure multi-party computation (MPC), so that
the comparison is evaluated jointly by two or three non-colluding servers
over secret-shared data~\cite{bozdemir2021privacy,fu2024ppa,fu2024fss}.
This keeps data owners offline and yields the most efficient results
reported to date. It rests, however, on a non-collusion assumption whose
violation reconstructs the original data, and it requires
server-to-server communication rounds that grow with the data size. The
non-collusion assumption also requires that the servers be run by
independent parties, for example a governmental institution and a
non-profit organization~\cite{bozdemir2021privacy}. A client that rents
computing from a single cloud provider cannot meet it on its own.

Neither approach keeps the comparison where the data is. We take a third
route and evaluate it inside the ciphertext. By approximating the sign
function with composite minimax polynomials~\cite{lee2021minimax}
evaluated under the Cheon--Kim--Kim--Song (CKKS)
scheme~\cite{cheon2017homomorphic}, core-point identification
and cluster propagation are completed on a single server. The client
uploads its encrypted data once and downloads the encrypted labels once;
between these two messages it sends and receives nothing, and neither the
client nor the server decrypts any intermediate value. In this sense, the
protocol is non-interactive. Existing privacy-preserving DBSCAN schemes
either involve the data owners in the
computation~\cite{amirbekyan2006privacy,kumar2007privacy,liu2012privacy,rahman2017towards,wang2022homomorphic}
or rely on multiple non-colluding
servers~\cite{bozdemir2021privacy,fu2024ppa,fu2024fss}; to the best of our
knowledge, no scheme avoids both.

Our main contributions are as follows.
\begin{itemize}
  \item We present the first non-interactive DBSCAN over encrypted data,
    in which every comparison is evaluated inside CKKS on a single server.
  \item We show that the encrypted result matches that of plaintext DBSCAN
    with the same border-point rule on six benchmark datasets and two
    real datasets from finance and healthcare.
  \item We give the client an algorithm that helps it set the number of
    label-propagation rounds before encryption: a grid-based estimate
    that uses only cell indices and cell counts and never compares a
    distance with the radius.
\end{itemize}